\chapter{Image Quality}

\section*{Patient focus}
This chapter explains the examination or clinical problem in accessible language. Patients should use it to prepare questions and understand the purpose, experience, benefits, and limitations of imaging; local instructions and the official report take priority.


\label{ch:image_quality}

Diagnostic accuracy relies on the quality of the acquired image. An image must have sufficient detail, contrast, and clarity to allow a radiologist to differentiate subtle pathologies from normal tissues. This chapter defines the core parameters of image quality—spatial, contrast, and temporal resolution—and reviews the factors that degrade them.

---

\section*{1. Core Parameters of Image Quality}

\subsection*{Spatial Resolution}
Spatial resolution is the ability of an imaging system to distinguish two high-contrast objects placed close together as separate entities. 
\textbullet \textbf{Measurement}: Quantified in line pairs per millimeter (lp/mm) or by the \textbf{Modulation Transfer Function (MTF)}.
\textbullet \textbf{Modality Comparison}: Mammography and planar X-ray have the highest spatial resolution ($\approx 15$ lp/mm), while Nuclear Medicine (PET/SPECT) has the lowest ($\approx 0.1$ lp/mm).
\textbullet \textbf{Limiting Factors}: Pixel size in digital detectors, focal spot size in X-ray tubes, and reconstruction slice thickness.

\subsection*{Contrast Resolution}
Contrast resolution is the ability to distinguish subtle differences in soft tissue density or signal intensity, even if the structures are adjacent.
\textbullet \textbf{Modality Comparison}: MRI has the highest contrast resolution, followed by CT. Planar X-ray has poor contrast resolution because all overlapping soft tissues are projected onto a single 2D plane.
\textbullet \textbf{Measurement}: Described by the \textbf{Contrast-to-Noise Ratio (CNR)}.

\subsection*{Temporal Resolution}
Temporal resolution is the ability to capture motion or rapid physiological changes without motion blur.
\textbullet \textbf{Modality Comparison}: Ultrasound has excellent temporal resolution, displaying real-time movement (up to 100 frames/second). Cardiac CT utilizes rapid gantry rotation times ($< 250$ ms) and ECG-gating to image the beating heart without blur.
\textbullet \textbf{Limiting Factors}: Scanner acquisition speed and detector response times.

---

\section*{2. Noise: The Signal-to-Noise Ratio (SNR)}
Image noise is the random, unwanted variation in pixel values that degrades image clarity, appearing as a "grainy" or "sandy" texture (mottle):

\textbullet \textbf{Quantum Mottle}: The most common source of noise in X-ray and CT. It is caused by an insufficient number of X-ray photons reaching the detector. If the dose (photon count) is too low, the image will appear highly grainy, obscuring subtle structures.
\textbullet \textbf{Overcoming Noise}: In X-ray/CT, noise is reduced by increasing the tube current-time product (mAs) or using iterative reconstruction. In MRI, noise is reduced by increasing the number of excitations (NEX) or utilizing larger voxel volumes, though this increases scan time or reduces spatial resolution.

---

\section*{3. Factors Degrading Image Quality}
\textbullet \textbf{Geometric Blurring (Penumbra)}: Caused by the finite size of the X-ray tube focal spot. It is minimized by keeping the patient as close to the detector as possible (minimizing object-to-detector distance).
\textbullet \textbf{Patient Motion}: Blurs anatomical boundaries. Minimized via patient coaching, physical immobilization, rapid acquisition sequences, or motion-correction algorithms.
\textbullet \textbf{Artifacts}: Systemic distortions (e.g. beam hardening in CT, chemical shift or metallic susceptibility in MRI, reverberation in ultrasound) that can mimic or obscure pathology.

---

\section*{4. The Imaging Triad: Quality vs. Safety vs. Time}
Every imaging acquisition requires balancing three competing factors:
1. \textbf{Diagnostic Quality}: Achieving high spatial and contrast resolution.
2. \textbf{Patient Safety / Dose}: In ionizing modalities (CT/X-ray), increasing image quality (lower noise) requires increasing the radiation dose, which must be balanced against the \textbf{ALARA (As Low As Reasonably Achievable)} safety principle.
3. \textbf{Acquisition Time}: In MRI, higher resolution or SNR requires longer scan times, which increases the likelihood of patient motion artifacts and reduces scanner throughput.

\clearpage
